\section{Introduction}
\label{sec:introduction}
Let $V=\{1,\ldots,n\}$, where $n\ge4$, and let
$d=(d_1,\ldots,d_n)\in\{0,\ldots,n-1\}^n$ be graphical.
Write $\Omega_d$ for the nonempty set of simple undirected graphs on $V$
whose labeled degree vector is $d$, and let $\pi_d$ be the uniform
probability measure on $\Omega_d$. A \emph{switch} removes two disjoint
edges and inserts a different perfect matching on their four vertices,
provided both new edges are absent. This preserves simplicity and
every vertex degree.

We use the following precise discrete-time convention. With probability
$1/2$ the chain stays at its current graph. Otherwise it chooses a
four-element set $S\subseteq V$ uniformly, and then chooses uniformly
one of the six ordered pairs $(F,F')$ of distinct perfect matchings of
the complete graph on $S$. If both edges of $F$ are present and both
edges of $F'$ are absent, it replaces $F$ by $F'$; otherwise it stays
put. Denote the transition kernel by $P_d$. For probability measures
on $\Omega_d$, set
\[
 \|\mu-\nu\|_{\mathrm{TV}}
 =\frac12\sum_{G\in\Omega_d}|\mu(G)-\nu(G)|,
\]
and define
\[
 t_{\mathrm{mix}}(d)
 =\min\left\{t\in\mathbb Z_{\ge0}:
 \max_{G\in\Omega_d}\|P_d^t(G,\cdot)-\pi_d\|_{\mathrm{TV}}
 \le\frac14\right\}.
\]

\begin{theorem}\label{thm:main}
For every $n\ge4$ and every graphical degree vector $d$,
\[
 t_{\mathrm{mix}}(d)\le 2n^8.
\]
If $|\Omega_d|>1$, the spectral gap of $P_d$ is at least
$[24n^2\binom n4]^{-1}$.
\end{theorem}

\begin{samepage}
\begin{corollary}[Exact uniform sampling]\label{cor:exact-sampling}
There is a uniform randomized algorithm that, given any finite graphical
labeled degree vector $d=(d_1,\ldots,d_n)$, uses unbiased random bits,
terminates almost surely, and outputs a graph chosen exactly uniformly
from the simple undirected graphs on $\{1,\ldots,n\}$ with labeled
degree vector $d$. Its expected bit running time is bounded by a fixed
polynomial in the binary encoding length of $d$.
\end{corollary}

The corollary includes vectors of length less than four, which are
handled directly. Section~\ref{sec:exact-sampling} proves it by
correcting the finite-time law of the chain. The time bound is in
expectation; no polynomial bound on every random execution is asserted.
\end{samepage}

\paragraph{History and significance.}
Uniform sampling with prescribed degrees provides a basic null model
for networks and a natural problem in approximate counting.

Kannan, Tetali, and Vempala initiated the switch-chain program, treated
regular bipartite sequences, and raised the possibility of broad
applicability~\cite{KTV99}. For simple undirected graphs, Cooper, Dyer,
and Greenhill established rapid mixing for regular sequences; their
corrigendum preserves the polynomial bound with a corrected
exponent~\cite{CDG07,CDG12}. Feder, Guetz, Mihail, and Saberi studied
early irregular families~\cite{FGMS06}, and Greenhill and Sfragara
proved rapid mixing under quantitative restrictions on the degrees,
correcting and extending Greenhill's earlier irregular-degree
argument~\cite{Greenhill15,GS18}. Amanatidis and Kleer developed a route through
strong stability~\cite{AK20}.

Stability conditions for degree-constrained sampling were developed by
Jerrum and Sinclair~\cite{JS90} and by Jerrum, McKay, and
Sinclair~\cite{JMS92}. Erd\H{o}s, Greenhill, Mezei, Mikl\'os,
Solt\'esz, and Soukup subsequently
proved rapid mixing for every P-stable family~\cite{Unified22}.
P-stability requires a polynomial bound, uniform over the family, on
the ratio between the number of realizations of nearby perturbed
sequences and that of the original sequence. Recent work of
Erd\H{o}s, Mikl\'os, and Soukup~\cite{EMS25} and of Erd\H{o}s, Lippner,
Nevo, and Soukup~\cite{ELNS25} gives further P-stable families from
regions in which every admissible degree vector is graphical.
These results explain substantial classes of degrees; the theorem
here requires only that the individual degree vector be graphical.
The modern Kannan--Tetali--Vempala conjecture predicts rapid mixing
for every realizable degree specification in each of the bipartite, directed, and
simple-undirected models; see \cite[Conjecture~1.1]{Unified22}.
Theorem~\ref{thm:main} resolves its simple-undirected formulation
affirmatively. Our off-diagonal transition probabilities agree with
the simple-undirected kernel displayed in
\cite[Equation~(3.1)]{Unified22}.

Fu, Qin, and Wang proved the bipartite formulation, equivalently the
fixed-margin binary-matrix problem~\cite{FQW26}. They compare the
swap chain with a two-row heat bath and reduce its spectral analysis
to a three-row conditional-projection inequality. We reprove the
equivalent variance inequality of their Theorem~5.1, retaining the
separation of label assignments and block counts and the adjacent-count
coupling developed in their Lemmas~5.2--5.4. A simple undirected graph
has a symmetric adjacency matrix, so an arbitrary binary-matrix
resampling does not preserve its state space. Our pair resamplings
preserve that symmetry by redistributing neighbors between two vertices.

These pair updates belong to the Curveball framework. For binary
matrices with fixed margins, Verhelst developed a multiple-swap
procedure~\cite{Verhelst08}. Strona, Nappo, Boccacci, Fattorini, and
San-Miguel-Ayanz later independently formulated a version called the
Curveball algorithm~\cite{Strona14,Strona16,CK18}; Carstens, Berger,
and Strona developed its undirected version~\cite{CBS18}.
They also observed that undirected trades on disjoint vertex pairs can
be dependent. The additional estimate here gives quantitative control
of those interactions: although an individual disjoint-pair term may
be negative, their full sum is positive semidefinite.

For prescribed degrees on arbitrary explicitly listed simple undirected
unweighted hosts, OpenAI~\cite[Corollary~10.1]{OpenAIPerfectMatchings2026}
gives a fully polynomial randomized approximation scheme for counting
labeled subgraphs with those degrees and, when feasible, a sampler within
total-variation distance $\tau$ of their uniform law. For rational
$0<\tau<1$, the sampling bit cost on every execution is polynomial in the
full input length and $\tau^{-1}$. These broader-host guarantees do not
assert mixing of a host-restricted switch chain: Theorem~\ref{thm:main}
concerns the complete host, while Corollary~\ref{cor:exact-sampling}
gives exact uniform sampling with expected polynomial bit cost.

\paragraph{Proof strategy.}
For an unordered vertex pair $a$, we resample which of its two endpoints
receives each neighbor adjacent to exactly one of them. All other
adjacencies and the number assigned to each endpoint remain fixed.
Each conditional space is therefore a uniform subset slice: all subsets
of one fixed size from the available neighbors. Let $E_a$ be its
averaging projection, $h_a=I-E_a$, and $H=\sum_a h_a$.
The central estimate is $H^2\succeq H$, where $\succeq$ denotes the
order of quadratic forms. Once the kernel of $H$ is shown to consist
of constants, comparison with individual switches gives the stated
gap and mixing time.

The reduction from $H$ to its restrictions on triples follows the
squared-generator method of Caputo~\cite[Section~1.1]{Caputo08}, also
used for binary matrices by Fu, Qin, and Wang~\cite[Lemma~4.1]{FQW26}.
The expansion groups terms according to whether the two indexing pairs
intersect. Caputo's disjoint terms are controlled by commutation
\cite[Equation~(1.11)]{Caputo08}.
Here the undirected trades need a separate estimate; this is where
the slice comparison and the disjoint-pair calculation enter.

First, all three pair resamplings on a vertex triple can be represented
by updates of a three-row binary matrix with fixed margins. When needed,
one extra column records the internal edges of the triple. We prove that the
variance of a sum of three row functions is at most twice the sum of
their variances. To see the structure behind this estimate, discard
columns of sum zero or three. Each remaining column chooses either
the row containing its unique one or the row containing its unique zero.
We first control these label assignments with their counts fixed,
then control the random counts themselves. The resulting variance
bound controls all overlapping pair resamplings.

Second, after fixing the data preserved by two disjoint pair
resamplings, the dependence between their conditional data has rank
at most one. Its negative contribution is bounded by projections onto
centered indicators that two neighbors receive the same endpoint.
A slice inequality bounds the sum of these projections by the switch
energy available on the same slice. Summing cancels the negative
contributions and completes the estimate for $H$. The three-row
variance inequality and the slice projection comparison are stated
separately from the graph argument, making their scope explicit.

\paragraph{Organization.}
Section~\ref{sec:setup} defines the pair resamplings and their conditional
projections. Section~\ref{row:section} proves the three-row inequality,
and Section~\ref{tri:section} applies it to overlapping updates.
Section~\ref{dis:section} identifies the errors in one disjoint-pair
interaction. Section~\ref{sl:section} controls their sum by a slice
comparison and proves positivity of the full disjoint contribution.
Section~\ref{mix:section} combines these estimates, proves connectivity,
and tracks the comparison and proposal constants through the mixing
bound. Section~\ref{sec:exact-sampling} then gives an exact uniform
sampler with expected polynomial bit cost. All proof ingredients are
supplied in the paper.
